\chapter*{Závěr}
\addcontentsline{toc}{chapter}{Závěr}
Cílem této bakalářské práce bylo navrhnout a implementovat mobilní aplikaci pro podporu lokálního prodeje produktů menších producentů.
Tento cíl byl splněn vytvořením aplikace, která umožňuje vyhledávání obchodů v~okolí uživatele, jejich filtrování podle různých kritérií, vytváření vlastních obchodů a práci s~uživatelskými profily a recenzemi.

V~rámci práce byl navržen architektonický koncept aplikace, který byl následně realizován pomocí moderních technologií a postupů.
Důraz byl kladen zejména na oddělení jednotlivých vrstev aplikace, přehlednost kódu a jeho dlouhodobou udržovatelnost.
Velmi přínosným prvkem se ukázalo zavedení vlastního typu \texttt{Domain\-Result} pro reprezentaci výsledků operací, inspirovaného monadickým přístupem, který umožnil konzistentní a robustní řešení chyb.
Celkově se zvolený návrh osvědčil při implementaci i postupném rozšiřování aplikace.

Součástí práce byla také implementace datové a doménové vrstvy, uživatelského rozhraní a bezpečnostního modelu.
Řešeny byly také specifické problémy, jako je například geografické vyhledávání, správa stavu uživatelského rozhraní a práce s~vedlejšími efekty nebo řízení přístupu k~datům.
Aplikace byla následně otestována pomocí jednotkových testů, manuálního testování i uživatelského testování, na jejichž základě došlo k~dalším úpravám a zlepšení použitelnosti.

Výsledkem práce je funkční aplikace, která je připravena pro používání reálnými uživateli a je dostupná v~obchodě Google Play.
Práce tak splnila stanovené cíle a přinesla praktický výstup ve formě reálně použitelné mobilní aplikace.

Dalším logickým krokem, který již přesahuje rámec této práce, je propagace aplikace mezi komunitami odpovídajícími cílovým uživatelům, k~čemuž je nyní aplikace připravena.